\documentclass[10pt]{article}

% Ћирилица
\usepackage[T2A]{fontenc}
\usepackage[utf8]{inputenc}
\usepackage[serbianc]{babel}

\usepackage[a4paper,top=1.5cm,bottom=1.5cm,left=1.8cm,right=1.8cm]{geometry}
\usepackage{amsmath}
\usepackage{multicol}
\usepackage{titlesec}
\titleformat*{\section}{\large\bfseries}
\titlespacing*{\section}{0pt}{6pt}{3pt}
\usepackage{enumitem}
\usepackage{tcolorbox}
\tcbuselibrary{breakable}
\usepackage[european]{circuitikz}
\usetikzlibrary{arrows.meta, babel}

\pagestyle{empty}
\setlength{\columnsep}{0.8cm}
\setlength{\parindent}{0pt}
\setlength{\parskip}{3pt}

% Оквир за кључне податке
\newtcolorbox{kljucno}{colback=blue!4, colframe=blue!45!black,
    boxrule=0.5pt, arc=1.5pt, left=6pt, right=6pt, top=3pt, bottom=3pt}

% Решени пример
\newtcolorbox[auto counter]{primer}{
    colback=white, colframe=black!55, colbacktitle=black!55, coltitle=white,
    fonttitle=\bfseries\small, title={Пример~\thetcbcounter},
    boxrule=0.5pt, arc=1.5pt, left=5pt, right=5pt, top=2pt, bottom=2pt,
    before skip=5pt, after skip=5pt}

% Задатак за вежбу (решење се даје само као резултат)
\newtcolorbox[auto counter]{zadatak}{
    colback=green!3, colframe=green!40!black, colbacktitle=green!40!black, coltitle=white,
    fonttitle=\bfseries\small, title={Задатак~\thetcbcounter},
    boxrule=0.5pt, arc=1.5pt, left=5pt, right=5pt, top=2pt, bottom=2pt,
    before skip=5pt, after skip=5pt}

% Цртеж проводника
\newcommand{\provodnik}{%
\begin{tikzpicture}[scale=0.8, >=Stealth, font=\small]
    \draw[thick, fill=orange!25] (0,0) ellipse (0.18 and 0.4);
    \draw[thick] (0,0.4) -- (3.6,0.4);
    \draw[thick] (0,-0.4) -- (3.6,-0.4);
    \draw[thick, fill=orange!15] (3.6,0) ellipse (0.18 and 0.4);
    \node at (0,0) {$S$};
    \draw[<->] (0,-0.75) -- (3.6,-0.75) node[midway, below] {$l$};
    \node at (1.8,0) {$\rho$};
\end{tikzpicture}}

% Цртеж плочастог кондензатора
\newcommand{\plocasti}{%
\begin{tikzpicture}[scale=0.8, >=Stealth, font=\small]
    \fill[blue!8] (0,0) rectangle (3,1);
    \draw[ultra thick] (0,0) -- (3,0);
    \draw[ultra thick] (0,1) -- (3,1);
    \draw (1.5,1) -- (1.5,1.4);
    \draw (1.5,0) -- (1.5,-0.4);
    \node at (1.5,0.5) {$\varepsilon$};
    \draw[<->] (3.3,0) -- (3.3,1) node[midway, right] {$d$};
    \node[above left] at (0.9,1) {$S$};
\end{tikzpicture}}

\begin{document}

\begin{center}
    {\Large\bfseries Отпорност, капацитивност и везе отпорника -- примери}
\end{center}
\vspace{-6pt}

%==================================================================
\section*{1. Префикси и степени броја 10}
%==================================================================

\begin{minipage}[t]{0.36\linewidth}
\vspace{0pt}
\centering\small
\renewcommand{\arraystretch}{1.1}
\begin{tabular}{|c|l|c|}
    \hline
    \textbf{Ознака} & \textbf{Назив} & \textbf{Степен}  \\ \hline
    T     & тера  & $10^{12}$   \\ \hline
    G     & гига  & $10^{9}$    \\ \hline
    M     & мега  & $10^{6}$    \\ \hline
    k     & кило  & $10^{3}$   \\ \hline
    --    & --    & $10^{0} = 1$ \\ \hline
    m     & мили  & $10^{-3}$  \\ \hline
    $\mu$ & микро & $10^{-6}$   \\ \hline
    n     & нано  & $10^{-9}$   \\ \hline
    p     & пико  & $10^{-12}$  \\ \hline
    f     & фемто & $10^{-15}$ \\ \hline
\end{tabular}
\end{minipage}\hfill
\begin{minipage}[t]{0.62\linewidth}
\vspace{0pt}
\begin{kljucno}
\textbf{Позитиван степен:} $10^{x}$ је јединица и иза ње $x$ нула:
$10^{3} = 1\,000$, \ $10^{6} = 1\,000\,000$.\\
\textbf{Негативан степен:} $10^{-x}$ је $0$, зарез, па јединица на
$x$-том месту иза зареза: $10^{-1} = 0{,}1$, \ $10^{-3} = 0{,}001$.\\
\textbf{Множење и дељење:} изложиоци се сабирају, односно одузимају:
\[ 10^{a}\cdot 10^{b} = 10^{a+b}, \qquad \frac{10^{a}}{10^{b}} = 10^{a-b} \]
\vspace{-12pt}
\end{kljucno}

\textbf{Претварање:} префикс заменимо његовим степеном:
$4{,}7~\text{k}\Omega = 4{,}7\cdot 10^{3}~\Omega = 4\,700~\Omega$, \;
$47~\text{nF} = 47\cdot 10^{-9}~\text{F} = 0{,}047~\mu\text{F}$, \;
$1~\text{mm}^2 = 10^{-6}~\text{m}^2$, \; $1~\text{cm}^2 = 10^{-4}~\text{m}^2$.
\end{minipage}

\vspace{4pt}
\begin{minipage}[t]{0.50\linewidth}
\vspace{0pt}
\centering\small
\renewcommand{\arraystretch}{1.1}
\begin{tabular}{|l|l|l|}
    \hline
    \textbf{Са префиксом} & \textbf{Са степеном} & \textbf{Као број} \\ \hline
    $3~\text{k}\Omega$      & $3\cdot 10^{3}~\Omega$       & $3\,000~\Omega$ \\ \hline
    $2{,}5~\text{M}\Omega$  & $2{,}5\cdot 10^{6}~\Omega$   & $2\,500\,000~\Omega$ \\ \hline
    $470~\text{mA}$         & $470\cdot 10^{-3}$~A         & $0{,}47$~A \\ \hline
    $22~\mu\text{F}$        & $22\cdot 10^{-6}$~F          & $0{,}000\,022$~F \\ \hline
    $15~\text{pF}$          & $15\cdot 10^{-12}$~F         & $0{,}000\,000\,000\,015$~F \\ \hline
\end{tabular}
\end{minipage}\hfill
\begin{minipage}[t]{0.47\linewidth}
\vspace{0pt}
\textbf{Од броја до префикса} иде обрнуто: тражимо степен 10 који
одговара неком префиксу:
\begin{align*}
0{,}005~\text{A} &= 5\cdot 10^{-3}~\text{A} = 5~\text{mA}\\
68\,000~\Omega &= 68\cdot 10^{3}~\Omega = 68~\text{k}\Omega\\
1\,200\,000~\text{Hz} &= 1{,}2\cdot 10^{6}~\text{Hz} = 1{,}2~\text{MHz}
\end{align*}
\end{minipage}

\begin{kljucno}
\textbf{Померање зареза:} суседни префикси (k, M, G или m, $\mu$, n, p)
разликују се $1000$ пута, тј. за \textbf{3 места} зареза.
Прелазак на \textbf{већи} префикс -- зарез иде \textbf{улево} (број се
смањује); на \textbf{мањи} префикс -- зарез иде \textbf{удесно} (број се
повећава):
$4\,700~\Omega = 4{,}7~\text{k}\Omega = 0{,}0047~\text{M}\Omega$, \quad
$0{,}22~\mu\text{F} = 220~\text{nF} = 220\,000~\text{pF}$.
\end{kljucno}

\textbf{Рачунање са величинама}
\vspace{-8pt}

\begin{multicols}{2}

\begin{primer}
\textbf{Множење.} Струја $I = 2$~mA тече кроз $R = 3~\text{k}\Omega$.
Колики је напон $U = I\cdot R$?
\[
U = 2\cdot 10^{-3}\cdot 3\cdot 10^{3} = 6\cdot 10^{-3+3} = 6\cdot 10^{0} = 6~\text{V}
\]
\end{primer}

\begin{primer}
\textbf{Дељење.} Напон $U = 12$~V је на отпорнику $R = 4~\text{k}\Omega$.
Колика је струја $I = U/R$?
\[
I = \frac{12}{4\cdot 10^{3}} = 3\cdot 10^{-3}~\text{A} = 3~\text{mA}
\]
\end{primer}

\columnbreak

\begin{primer}
\textbf{Сабирање.} $2{,}2~\text{k}\Omega + 470~\Omega = ?$
Сабирају се само величине \textbf{са истим префиксом}, зато прво све
пребацимо у оме:
$2\,200~\Omega + 470~\Omega = 2\,670~\Omega = 2{,}67~\text{k}\Omega$
\end{primer}

\begin{primer}
\textbf{Одузимање.} $1~\mu\text{F} - 470~\text{nF} = ?$
$1~\mu\text{F} = 1\,000~\text{nF}$, па је
$1\,000~\text{nF} - 470~\text{nF} = 530~\text{nF} = 0{,}53~\mu\text{F}$
\end{primer}

\end{multicols}
\vspace{-10pt}

\begin{zadatak}
Попуни: а) $3{,}3~\text{k}\Omega = \_\_\cdot 10^{3}~\Omega = \_\_~\Omega$; \;
б) $0{,}000\,47~\text{F} = \_\_\cdot 10^{-6}~\text{F} = \_\_~\mu\text{F}$; \;
в) $10^{5}\cdot 10^{-2} = ?$; \; г) $\dfrac{10^{-2}}{10^{3}} = ?$; \;
д) $5~\text{mA}\cdot 2~\text{k}\Omega = ?$
\hfill\textit{[а) 3,3 и 3\,300; б) 470 и 470; в) $10^{3}$; г) $10^{-5}$; д) 10~V]}
\end{zadatak}

%==================================================================
\section*{2. Проценат}
%==================================================================

\textbf{Проценат} (\%) значи „од сто”: $1\,\% = \frac{1}{100} = 0{,}01$.
\quad \textbf{$p$\,\% од броја $A$} је $A\cdot\dfrac{p}{100}$
\quad (нпр. $20\,\%$ од $50~\Omega$ је $50\cdot 0{,}2 = 10~\Omega$).
\quad \textbf{Промена у процентима:} $\dfrac{\text{нова} - \text{стара}}{\text{стара}}\cdot 100\,\%$.

У електротехници се проценат најчешће среће као \textbf{толеранција}: отпорник
$100~\Omega \pm 5\,\%$ може имати било коју вредност од $95$ до $105~\Omega$.

\newpage

%==================================================================
\section*{3. Лаки примери}
%==================================================================

\begin{multicols}{2}

\begin{primer}
Проводник: $\rho = 0{,}02~\Omega\,\text{mm}^2/\text{m}$, $l = 100$~m,
$S = 2~\text{mm}^2$. Колика је отпорност?

$R = \rho\,\dfrac{l}{S} = 0{,}02\cdot\dfrac{100}{2} = 0{,}02\cdot 50 = 1~\Omega$
\end{primer}

\begin{primer}
Кондензатор: $\varepsilon = 2$~F/m, $S = 3~\text{m}^2$, $d = 0{,}5$~m.
Колика је капацитивност?

$C = \varepsilon\,\dfrac{S}{d} = 2\cdot\dfrac{3}{0{,}5} = 2\cdot 6 = 12$~F
\end{primer}

\begin{primer}
Редна веза: $R_1 = 2~\Omega$, $R_2 = 3~\Omega$, $R_3 = 5~\Omega$.

\begin{center}
\begin{circuitikz}[scale=0.65, transform shape]
    \draw (0,0) to[R, l=$R_1$, o-] (2,0) to[R, l=$R_2$] (4,0) to[R, l=$R_3$, -o] (6,0);
\end{circuitikz}
\end{center}
$R_e = 2 + 3 + 5 = 10~\Omega$
\end{primer}

\begin{primer}
Паралелна веза: $R_1 = 6~\Omega$, $R_2 = 3~\Omega$.

\begin{center}
\begin{circuitikz}[scale=0.65, transform shape]
    \draw (0,0) to[short, o-*] (1,0);
    \draw (1,0) -- (1,0.8) to[R, l=$R_1$] (4,0.8) -- (4,0);
    \draw (1,0) -- (1,-0.8) to[R, l_=$R_2$] (4,-0.8) -- (4,0);
    \draw (4,0) to[short, *-o] (5,0);
\end{circuitikz}
\end{center}
$R_e = \dfrac{6\cdot 3}{6 + 3} = \dfrac{18}{9} = 2~\Omega$
\end{primer}

\begin{zadatak}
Израчунај $R$ за $\rho = 0{,}02~\Omega\,\text{mm}^2/\text{m}$, $l = 50$~m, $S = 1~\text{mm}^2$.
\hfill\textit{[$R = 1~\Omega$]}
\end{zadatak}

\begin{zadatak}
Израчунај $C$ за $\varepsilon = 3$~F/m, $S = 4~\text{m}^2$, $d = 2$~m.
\hfill\textit{[$C = 6$~F]}
\end{zadatak}

\begin{zadatak}
Колика је $R_e$: а) редно $1{,}5~\Omega$ и $2{,}5~\Omega$; б) паралелно $4~\Omega$ и $4~\Omega$?
\hfill\textit{[а) $4~\Omega$; б) $2~\Omega$]}
\end{zadatak}

\end{multicols}

%==================================================================
\section*{4. Примери са префиксима}
%==================================================================

\begin{multicols}{2}

\begin{primer}
Алуминијумски вод: $\rho = 2{,}8\cdot 10^{-8}~\Omega\,\text{m}$, $l = 2$~km,
$S = 4~\text{mm}^2$. Колика је отпорност?

Све у основне јединице: $l = 2\cdot 10^{3}$~m, $S = 4\cdot 10^{-6}~\text{m}^2$.
\begin{align*}
R &= 2{,}8\cdot 10^{-8}\cdot\frac{2\cdot 10^{3}}{4\cdot 10^{-6}}
   = 2{,}8\cdot 0{,}5\cdot 10^{-8+3+6}\\
  &= 1{,}4\cdot 10^{1} = 14~\Omega
\end{align*}
\end{primer}

\begin{primer}
Плочасти кондензатор: $\varepsilon_r = 4$, $S = 100~\text{cm}^2$, $d = 0{,}1$~mm.
Колика је капацитивност?

$S = 100\cdot 10^{-4} = 10^{-2}~\text{m}^2$, \; $d = 0{,}1\cdot 10^{-3} = 10^{-4}$~m.
\begin{align*}
C &= 8{,}85\cdot 10^{-12}\cdot 4\cdot\frac{10^{-2}}{10^{-4}}
   = 35{,}4\cdot 10^{-12}\cdot 10^{2}\\
  &= 3{,}54\cdot 10^{-9}~\text{F} = 3{,}54~\text{nF}
\end{align*}
\end{primer}

\begin{primer}
Редна веза: $R_1 = 2{,}2~\text{k}\Omega$, $R_2 = 470~\Omega$, $R_3 = 0{,}33~\text{k}\Omega$.

Све у омима: $R_e = 2\,200 + 470 + 330 = 3\,000~\Omega = 3~\text{k}\Omega$
\end{primer}

\begin{primer}
Отпорник $4{,}7~\text{k}\Omega \pm 10\,\%$. У ком опсегу је стварна отпорност?

$10\,\%$ од $4{,}7~\text{k}\Omega$ је $0{,}47~\text{k}\Omega$, па је
$R$ од $4{,}23$ до $5{,}17~\text{k}\Omega$.
\end{primer}

\begin{zadatak}
Бакарна жица: $\rho = 1{,}75\cdot 10^{-8}~\Omega\,\text{m}$, $l = 250$~m,
$S = 0{,}5~\text{mm}^2$. Израчунај $R$.
\hfill\textit{[$R = 8{,}75~\Omega$]}
\end{zadatak}

\begin{zadatak}
Кондензатор: $\varepsilon_r = 5$, $S = 200~\text{cm}^2$, $d = 0{,}2$~mm. Израчунај $C$.
\hfill\textit{[$C \approx 4{,}43$~nF]}
\end{zadatak}

\begin{zadatak}
$R_1 = 1{,}5~\text{k}\Omega$ везан је редно са паралелном везом
$R_2 = 3~\text{k}\Omega$ и $R_3 = 6~\text{k}\Omega$. Колика је $R_e$?
\hfill\textit{[$R_e = 3{,}5~\text{k}\Omega$]}
\vspace{-6pt}
\begin{center}
\begin{circuitikz}[scale=0.6, transform shape]
    \draw (0,0) to[R, l=$R_1$, o-*] (3,0);
    \draw (3,0) -- (3,0.8) to[R, l=$R_2$] (6,0.8) -- (6,0);
    \draw (3,0) -- (3,-0.8) to[R, l_=$R_3$] (6,-0.8) -- (6,0);
    \draw (6,0) to[short, *-o] (7,0);
\end{circuitikz}
\end{center}
\vspace{-8pt}
\end{zadatak}

\begin{zadatak}
Претвори: а) $0{,}047~\mu\text{F}$ у nF; б) $2{,}2~\text{M}\Omega$ у $\Omega$;
в) $15~\text{mA}$ у A. Колико је $15\,\%$ од $2{,}4~\text{k}\Omega$?
\hfill\textit{[а) $47$~nF; б) $2\,200\,000~\Omega$; в) $0{,}015$~A; $360~\Omega$]}
\end{zadatak}

\end{multicols}

\end{document}